\section{The surviving norm-character obstruction}\label{sec:obstruction}

Theorem~\ref{thm:jet-comparison} retains a line whose restriction to \(G_F\)
is the norm character in the upper two cohomological degrees. We now show
that this line cannot be a periodicity
power over the generic deformation field, and then use the actual finite
abutment filtration to contradict Proposition~\ref{prop:constituent-test}.

\subsection{The norm line over \texorpdfstring{\(Q_x\)}{Qx}}

Recall that \(G_F=\mathcal O_F^\times\) is the full unit group of the
unramified quadratic extension \(F/\mathbb Q_3\), acting semilinearly on
\(Q_x=k((x))\). Its residue norm is the surjective character
\[
 \delta:G_F\longrightarrow\mathbb F_9^\times
   \xrightarrow{\,N_{\mathbb F_9/\mathbb F_3}\,}\mathbb F_3^\times
   \subset k^\times.
\]
In particular \(\delta\) is nontrivial and \(\delta^2=1\). The following
lemma uses the whole group, including its principal units. Restricting only
to the finite residue-field subgroup would not yield this conclusion.

\begin{lemma}[A character not supplied by periodicity]
\label{lem:generic-character}
The fixed field \(Q_x^{G_F}\) is \(k\). Moreover,
\[
 Q_x(\delta)\not\simeq\overline\omega^{\otimes b}
             \qquad\text{for every }b\in\mathbb Z
\]
as semilinear \(G_F\)-representations.
\end{lemma}

\begin{proof}
We first show that the action of \(G_F\) on \(Q_x\) has infinite image.
Let \(\mathcal F\) be the characteristic-\(p\) universal height-two formal
group over \(\mathcal O_x\). If \(g\in G_F\) acts trivially on
\(\mathcal O_x\), its marking-change isomorphism is an automorphism of
\(\mathcal F\) over that ring. Over \(Q_x\), the first Hasse coefficient is
invertible, so \(\mathcal F\) has height one. After extending \(Q_x\) to an
algebraic closure, the height classification identifies it with the
multiplicative formal group \citep[Theorem~IV]{Lazard55}.

The automorphisms of the multiplicative formal group are precisely the
intrinsic scalar automorphisms \([c]\), \(c\in\Zp^\times\). Indeed,
restriction to its finite \(p^r\)-torsion group schemes identifies an
automorphism with a compatible system in
\(\lim_r(\mathbb Z/p^r)^\times=\Zp^\times\): the character group of
\(\mu_{p^r}\) is \(\mathbb Z/p^r\), and the ideals \((T^{p^r})\) are cofinal
in the formal-coordinate topology. Thus these restrictions determine one
scalar automorphism of the entire formal group.

Every \([c]\) is already defined on \(\mathcal F\) over \(\mathcal O_x\).
For example, integer approximations to \(c\) modulo \(p^r\) give compatible
formal power series because the order of \([p^r](T)\) tends to infinity.
The injection of \(\mathcal O_x\) into the algebraic closure of \(Q_x\)
therefore shows, coefficient by coefficient, that the marking-change
automorphism equals \([c]\) over \(\mathcal O_x\). Specializing at \(x=0\)
shows that \(g\) is scalar on the Honda group, or that \(g^{-1}\) is scalar
under the opposite action convention. Thus the kernel of the action on
\(Q_x\) is contained in the central subgroup \(\Zp^\times\).
The quotient \(\mathcal O_F^\times/\Zp^\times\) is infinite: on principal
units, the \(3\)-adic logarithm identifies the relevant ranks with those
of \(3\mathcal O_F\) and \(3\Zp\), namely two and one. The action on \(Q_x\)
has infinite image.

Every element of \(G_F\) fixes \(k\) and preserves the \(x\)-adic valuation.
Suppose \(f\in Q_x^{G_F}\) is not in \(k\). If its valuation is negative,
replace it by \(f^{-1}\); if its valuation is zero, subtract its residue in
\(k\). We obtain a nonzero invariant \(z\) with strictly positive valuation
\(d=v_x(z)\). The substitution \(k[[Z]]\to k[[x]]\), \(Z\mapsto z\), makes
\(k[[x]]\) finite free of rank \(d\) over \(k[[z]]\), by the one-variable
Weierstrass theorem. Hence \(Q_x/k((z))\) is a finite field extension.
Continuity implies that every element fixing \(z\) fixes \(k((z))\), so the
image of \(G_F\) lies in the finite automorphism group of this extension.
This contradicts its infinite image. Separability is not needed: the
automorphism group of any finite field extension has cardinality at most
its degree. We have proved
\begin{equation}\label{eq:generic-fixed-field}
                             Q_x^{G_F}=k.
\end{equation}

Next consider the first Hasse invariant. For an even-periodic theory the
cotangent line of its formal group identifies with \(\pi_2\): restriction
of the augmentation ideal to \(\mathbb{CP}^1\) identifies its quotient by
its square with \(\widetilde E_2^0(\mathbb{CP}^1)=\pi_2E_2\). Thus
\(\Omega=\omega/p\) is the cotangent line of \(\mathcal F\). In a coordinate
\(T\), let \(a_1\) be the coefficient of \(T^p\) in \([p]_{\mathcal F}(T)\).
Under \(T'=cT+O(T^2)\), that coefficient becomes \(c^{1-p}a_1\), while
\(dT'=c\,dT\) at the identity. Therefore
\[
                    h=a_1(dT)^{\otimes(p-1)}
\]
is an intrinsic, hence \(G_F\)-invariant, section of
\(\Omega^{\otimes(p-1)}=\Omega^{\otimes2}\). The universal height-two
parameter is its simple zero: the transported Lubin--Tate normalization
in Section~\ref{sec:setup} gives \(a_1=x\) in the chosen coordinate.
In any integral frame \(e\) of \(\Omega\), write
\begin{equation}\label{eq:hasse-valuation}
                    h=h_e e^{\otimes2},\qquad v_x(h_e)=1.
\end{equation}
Thus \(h\) is an invariant basis of \(\overline\omega^{\otimes2}\) over
\(Q_x\).

Suppose now that \(Q_x(\delta)\simeq\overline\omega^{\otimes b}\). The image
of the distinguished basis of \(Q_x(\delta)\) is a basis
\(v=f e^{\otimes b}\) with \(f\in Q_x^\times\) and
\(g(v)=\delta(g)v\). Since \(\delta^2=1\), the ratio
\[
               \frac{v^{\otimes2}}{h^{\otimes b}}
                     =\frac{f^2}{h_e^b}
\]
is a nonzero invariant function. By \eqref{eq:generic-fixed-field}, it lies
in \(k^\times\). Taking valuations in \eqref{eq:hasse-valuation} gives the
exact equality
\begin{equation}\label{eq:character-parity}
                              2v_x(f)=b.
\end{equation}
Hence \(b\) is even. But \(h^{\otimes b/2}\) is then an invariant basis of
\(\overline\omega^{\otimes b}\), so this line is trivial.

Finally, \(Q_x(\delta)\) is not trivial. Otherwise there would be
\(u\in Q_x^\times\) with \(g(u)=\delta(g)u\) for every \(g\in G_F\); the
inverse convention is the same because \(\delta^{-1}=\delta\). Then
\(u^2\in k^\times\) by \eqref{eq:generic-fixed-field}, so \(u\) is algebraic
over \(k\). The finite field \(k\) is relatively algebraically closed in
\(k((x))\). To see this, an element algebraic over \(k\) satisfies
\(u^{|k|^r}=u\) for some \(r>0\). If nonzero it has valuation zero;
subtract its residue and compare positive valuations in the same equation
to see that the difference is zero. Thus \(u\in k^\times\), contradicting
the nontriviality of \(\delta\). This proves the lemma.
\end{proof}

The fixed-field part of the proof is what prevents a nonconstant function
of \(x\) from disguising the norm character as a periodicity twist. The
remaining step is to ensure that the line cannot disappear before reaching
homotopy.

\subsection{The actual homotopy filtration}

We first record the homotopy calculation independently of the
finite-assembly question. Recall \(Z_k=L_{K(3)}S_k\), and write
\(\mathcal V_d^\Delta(Z_k)\) for the diagonal constant-field summand of
the test \(\mathcal V_d(Z_k)\) from \eqref{eq:generic-test}.
Its scalar field is \(Q_x\), and we restrict its action to \(G_F\).

\begin{proposition}[The generic diagonal homotopy]
\label{prop:generic-homotopy}
The diagonal spectral sequence of Proposition~\ref{prop:descent-test},
after extension of scalars to \(Q_x\), collapses at its second page.
For every \(r\in\mathbb Z\), its actual abutment filtration gives
\(G_F\)-equivariant short exact sequences
\begin{align}
 0&\longrightarrow
 Q_x(\delta)\otimes_{Q_x}\overline\omega^{\otimes(r+2)}
 \longrightarrow\mathcal V_{2r}^\Delta(Z_k)
 \longrightarrow\overline\omega^{\otimes r}\longrightarrow0,
 \label{eq:generic-even-filtration}\\
 0&\longrightarrow
 Q_x(\delta)\otimes_{Q_x}\overline\omega^{\otimes(r+1)}
 \longrightarrow\mathcal V_{2r-1}^\Delta(Z_k)
 \longrightarrow\overline\omega^{\otimes r}\longrightarrow0.
 \label{eq:generic-odd-filtration}
\end{align}
In particular, each diagonal summand \(\mathcal V_d^\Delta(Z_k)\) has
dimension two over \(Q_x\), and \(\mathcal V_{-3}^\Delta(Z_k)\) contains
a subrepresentation isomorphic to \(Q_x(\delta)\).
\end{proposition}

\begin{proof}

Apply Proposition~\ref{prop:descent-test} to
\(Y/p=\mathcal T(Z_k)\), take the diagonal constant-field summand, and
localize homotopy groups from \(\mathcal O_x\) to \(Q_x\). Restrict the
actions to \(G_F\). Theorem~\ref{thm:jet-comparison} gives the complete list
of nonzero terms on the second page:
\begin{equation}\label{eq:obstruction-page}
\begin{array}{c|c}
 \text{cohomological degree }s & \mathsf E_2^{s,2j},\quad j\in\mathbb Z
 \\\hline
 0,1 & \overline\omega^{\otimes j}\\[2pt]
 3,4 & Q_x(\delta)\otimes_{Q_x}\overline\omega^{\otimes j}.
\end{array}
\end{equation}
All odd internal degrees vanish. Here \(\mathsf E_r\) denotes the diagonal,
localized spectral sequence. Its differentials have bidegree
\[
 d_r:\mathsf E_r^{s,t}\longrightarrow\mathsf E_r^{s+r,t+r-1}.
\]
An even \(r\) changes internal parity, so its differential is zero. For odd
\(r\geq3\), the four cohomological degrees in \eqref{eq:obstruction-page}
leave only \(r=3\), from \(s=0\) to \(s=3\) or from \(s=1\) to \(s=4\).
Every such differential would have the form
\begin{equation}\label{eq:forbidden-differential}
 \overline\omega^{\otimes j}\xrightarrow{\ d_3\ }
       Q_x(\delta)\otimes_{Q_x}\overline\omega^{\otimes(j+1)}.
\end{equation}
It is a \(Q_x\)-linear, \(G_F\)-equivariant map. A nonzero map between these
one-dimensional representations is an isomorphism; after tensoring with
the inverse periodicity line, it would identify \(Q_x(\delta)\) with a
periodicity power. Lemma~\ref{lem:generic-character} excludes that
possibility. Both possible \(d_3\) families therefore vanish, and no later
differential is possible.

The diagonal idempotent commutes with the actual finite filtration
supplied by Proposition~\ref{prop:descent-test}, and localization from
\(\mathcal O_x\) to \(Q_x\) is exact. Thus
\[
 F^s\mathcal V_d^\Delta(Z_k)/F^{s+1}\mathcal V_d^\Delta(Z_k)
       \simeq\mathsf E_\infty^{s,d+s}.
\]
For \(d=2r\), the only nonzero quotients occur at \(s=0,4\);
they are \(\overline\omega^{\otimes r}\) and
\(Q_x(\delta)\otimes\overline\omega^{\otimes(r+2)}\), respectively.
For \(d=2r-1\), they occur at \(s=1,3\), with lines
\(\overline\omega^{\otimes r}\) and
\(Q_x(\delta)\otimes\overline\omega^{\otimes(r+1)}\).
The filtration starts with the whole homotopy group and ends at zero,
so the higher-filtration line is a subrepresentation and the lower one
is its quotient. This proves the two exact sequences and the dimension
assertion. Taking \(r=-1\) in \eqref{eq:generic-odd-filtration} gives
\[
 0\longrightarrow Q_x(\delta)
 \longrightarrow\mathcal V_{-3}^\Delta(Z_k)
 \longrightarrow\overline\omega^{\otimes(-1)}\longrightarrow0.
\]
\end{proof}

These sequences describe the two nonzero graded pieces and their order
in actual homotopy; they do not assert an equivariant splitting. The
norm line can now be compared directly with the finite-assembly detector.

\begin{proof}[Proof of Theorem~\ref{thm:main}]
Suppose, to the contrary, that
\[
          L_2L_{K(3)}S\in\Thick\{L_0S,L_1S,L_2S\}
\]
in the category specified in the Theorem. Forgetting to spectra and
applying exact \(K(2)\)-localization sends \(L_0S\) and \(L_1S\) to zero,
and sends \(L_2S\) to \(\mathbf 1_2\). Moreover,
\(L_{K(2)}L_2L_{K(3)}S\simeq L_{K(2)}L_{K(3)}S\). These are the
Bousfield-class relations in \citet[Theorem~2.1(d),(i)]{Ravenel84}, together
with the smashing property of \(E(i)\)-localization
\citep[Theorem~5.1 and Proposition~5.3]{HS99}. Thus
\(L_{K(2)}L_{K(3)}S\) lies in the ordinary thick subcategory generated by
\(\mathbf 1_2\).

By \eqref{eq:spherical-constants}, the underlying \(S\)-module \(S_k\) is
finite free. Consequently \(Z_k=L_{K(3)}S_k\), as an underlying spectrum,
is a finite direct sum of copies of \(L_{K(3)}S\). The same thick-generation
conclusion holds for \(L_{K(2)}Z_k\). No Galois-equivariant choice of basis
of \(S_k\) is required: the detector action comes from its \(E_2\) factor,
so every underlying spectrum map induces an equivariant detector map.
Proposition~\ref{prop:constituent-test} implies that
\begin{equation}\label{eq:required-constituents}
 \mathcal V_d(Z_k)\text{ has only constituents }
       \overline\omega^{\otimes b},\quad b\in\mathbb Z,
       \quad\text{for every }d\in\mathbb Z.
\end{equation}
But Proposition~\ref{prop:generic-homotopy} supplies the simple
subrepresentation \(Q_x(\delta)\) of the diagonal summand of
\(\mathcal V_{-3}(Z_k)\), hence a simple constituent of that whole test.
Lemma~\ref{lem:generic-character} excludes it from
\eqref{eq:required-constituents}. This contradiction proves the Theorem.
\end{proof}
